% Generated draft: requires mathematical and source review.
\documentclass[11pt]{article}
\usepackage[margin=25mm]{geometry}
\usepackage{fontspec}
\setmainfont{Latin Modern Roman}
\usepackage{amsmath,amssymb,mathtools,unicode-math}
\setmathfont{Latin Modern Math}
\usepackage{enumitem,needspace,xurl,hyperref}
\setlength{\parindent}{0pt}
\setlength{\parskip}{5pt}
\setlength{\emergencystretch}{3em}
\newcommand{\N}{\mathbb{N}}
\newcommand{\Z}{\mathbb{Z}}
\newcommand{\Q}{\mathbb{Q}}
\newcommand{\R}{\mathbb{R}}
\newcommand{\C}{\mathbb{C}}
\newcommand{\F}{\mathbb{F}}
\newcommand{\A}{\mathbb{A}}
\newcommand{\PP}{\mathbb{P}}
\newcommand{\E}{\mathbb{E}}
\newcommand{\eps}{\varepsilon}
\newcommand{\dd}{\,\mathrm{d}}
\begin{document}
\section*{UnsolvedMath ID: 30002361}
\textbf{Source problem code:} catalogue entry
\subsection{Definitions and mathematical statement}
Write \(\R/\Z\) for the additive quotient group of real numbers by integers. Retain \(KR^0\) and \(KR_{\R/\Z}^{-1}\) exactly as the source's notation for the proposed source and target functors. The supplied record does not define these functors, the meaning or hypotheses of \(R\), or the coefficient construction indicated by the subscript \(\R/\Z\); no identification with a particular version of \(K\)-theory is made here.

\textbf{Lott's secondary-index conjecture, as supplied.} There exists a natural transformation
\[
\eta\colon KR^0\longrightarrow KR_{\R/\Z}^{-1}
\]
that identifies analytic push-forward with the Becker--Gottlieb--Dold transfer.

The intended compatibility can be expressed schematically as follows. Write \(\mathcal F(X)=KR^0(X)\) and \(\mathcal G(X)=KR_{\R/\Z}^{-1}(X)\). For every map \(p\colon E\to B\) admitted by the intended theory, if the relevant operators have types
\[
p_*^{\mathrm{an}}\colon\mathcal F(E)\to\mathcal F(B),
\qquad
\operatorname{tr}_p\colon\mathcal G(E)\to\mathcal G(B),
\]
then the required identity is
\[
\eta_B\bigl(p_*^{\mathrm{an}}(x)\bigr)
=
\operatorname{tr}_p\bigl(\eta_E(x)\bigr)
\qquad\text{for every }x\in\mathcal F(E).
\]
Here \(p_*^{\mathrm{an}}\) denotes the intended analytic push-forward and \(\operatorname{tr}_p\) the intended Becker--Gottlieb--Dold transfer. This formula records the requested correspondence; its displayed operator types are conditional and are not additional hypotheses supplied by the record. Naturality means that the components \(\eta_X\) commute with the functorial maps in the intended category.

\textbf{Unresolved scope.} The record does not specify that category, its objects and morphisms, the admissible maps \(p\), the geometric data required for analytic push-forward, or any degree shifts and boundary conditions. Thus it does not determine a fully specified, independently interpretable conjecture. These details must be recovered from the cited sources before the schematic formulation can be replaced by a precise quantified statement.

\textbf{Supplied status qualification.} The dated curation labels the problem partially solved and reports that the required natural transformation has been constructed, while compatibility with the Becker--Gottlieb--Dold transfer remains unproved. This is a supplied literature assessment, not an independently verified mathematical or current-status claim.
\subsection{Short English statement}
The source attributes to Lott a conjecture asking for a natural transformation from \(KR^0\) to \(KR_{\R/\Z}^{-1}\) that carries analytic push-forward to the Becker--Gottlieb--Dold transfer. Supplied dated curation reports construction of the transformation but leaves transfer compatibility unproved. The functors, admissible maps, and geometric hypotheses remain unspecified in the record, so the statement requires source consultation and human mathematical review.
\subsection{Sources}
\begin{itemize}
\item Local UnsolvedMath registry, \texttt{data/unsolvedmath/problems.json}, supplied complete source record for permanent ID \(30002361\), source problem code \texttt{catalogue entry}, entitled \emph{Lott Secondary Index and Transfer Compatibility}. The record attributes the conjecture to Lott, gives proposed year 2013, and supplies no separate proposer identification. Its \texttt{corrected\_verified} statement label and claim of unchanged mathematical content are registry assertions, not verification performed for this draft.
\item Supplied report citation: \emph{Algebraic K-theory and Motivic Cohomology (2014)}, Oberwolfach Reports, pp.\ 1861--1913, DOI \texttt{10.4171/owr/2013/32}. \url{https://doi.org/10.4171/owr/2013/32}. The background describes these reports as records of mathematical workshops and their open-problem discussions. The citation's 2014 designation differs from the proposed year and year appearing in the DOI; the supplied information does not resolve this dating discrepancy.
\item Supplied literature source: \emph{Oberwolfach report on differential algebraic K-theory}. \url{https://ems.press/content/serial-article-files/46461?nt=1}.
\item Supplied literature source: Bunke and Gepner, \emph{differential K-theory transfer}. \url{https://arxiv.org/abs/1306.0247}.
\item Supplied literature source: Bunke and Gepner, \emph{Memoirs AMS}. \url{https://bookstore.ams.org/MEMO/269/1316}.
\item Literature assessment dated \texttt{2026-08-21}: the supplied curation reports that the natural transformation required by Lott's conjecture is constructed, but compatibility of analytic push-forward with the Becker--Gottlieb--Dold transfer remains unproved. This assessment is supplied dated curation derived from the OWR material, not independently verified here. Its time-sensitive status and cited claims require independent verification.
\end{itemize}
\end{document}
